%% kapa: exam=1 round=35 year=2024 period=1 form=A subject=economics subject_ko=경제학원론 seq=6 number=46
%% source: Q-Net 감정평가사 기출문제 (cst00309) artlSeq=5213229 file=raw/1008/5213229/1교시 문제지 원본.pdf
%% conversion: kordoc 4.24.51 -> tools/gen1.py (auto); review=proofread-sonnet; flags=none
\begin{question}{46}
완전경쟁시장에서 모든 기업이 이윤을 극대화하고 있는 산업 A는 비용곡선이 $C(Q)=2+\frac{Q^{2}}{2}$인 100개의 기업과 $C(Q)=\frac{Q^{2}}{10}$인 60개의 기업으로 구성되어 있다. 신규 기업의 진입이 없을 때, 가격이 2보다 큰 경우 산업 A의 공급곡선은? (단, $Q$는 생산량이다.)
\begin{choices}
  \item $Q=200P$
  \item $Q=300P$
  \item $Q=400P$
  \item $Q=415P$
  \item $Q=435P$
\end{choices}
\end{question}
\answer{3}
